% =====================================================================
%  DDDP-OTD: Parallel Multi-Cut Benders over Disjoint Temporal Windows
%  Drop-in section for a top-tier IEEE Transactions paper.
%  Requires:  amsmath, amssymb, algorithm, algpseudocode, bm
% =====================================================================

\section{Parallel Dual Dynamic Decomposition over Time (DDDP--OTD)}
\label{sec:dddp_otd}

The single-shot multi-period OPF over a horizon
$\mathcal{T}=\{1,\dots,T\}$ couples all timesteps through the battery
state-of-charge (SOC) dynamics
\begin{equation}
\label{eq:soc_global}
B_{t,j} \;=\; B_{t-1,j}\;+\;\eta_c\,P^{c}_{t,j}\;-\;\eta_d^{-1}\,P^{d}_{t,j},
\qquad \forall\,t\!\in\!\mathcal{T},\,j\!\in\!\mathcal{B},
\end{equation}
which makes the monolithic problem solve-time scale super-linearly with
$T$. We therefore decompose~\eqref{eq:soc_global} \emph{temporally} into a
sequence of smaller window problems, coupled only by the SOC handed
across window boundaries, and recover the centralized optimum through a
parallel multi-cut Benders scheme that we call \emph{Dual Dynamic
Decomposition over Time} (DDDP--OTD).

% ---------------------------------------------------------------
\subsection{Window Construction and Boundary Variables}
\label{ssec:windows}

Partition $\mathcal{T}$ into $P$ \emph{disjoint} contiguous windows
$\{\mathcal{W}_i\}_{i=1}^{P}$, where window $i$ covers the index range
$\mathcal{W}_i = \{w^{s}_i,\dots,w^{e}_i\}$ with
$w^{s}_1 = 1$, $w^{e}_P = T$, and
$w^{s}_{i+1} = w^{e}_i + 1$. In contrast to the overlapping Schwarz
variants commonly adopted for temporal OPF \cite{XX}, the windows here
carry \emph{no} look-ahead region: the right-side influence of future
stages is represented \emph{exactly} by Benders cuts rather than
approximately by an overlap penalty.

Let $\mathcal{B}$ denote the set of battery buses. The temporal
coupling between two adjacent windows is mediated by a single boundary
vector
\begin{equation}
\label{eq:bnd_def}
\bm{x}_i \;\triangleq\; \big(B_{w^{e}_i,\,j}\big)_{j\in\mathcal{B}}
\;\in\;\mathbb{R}^{|\mathcal{B}|},
\end{equation}
i.e.\ the terminal SOC produced by window $i$. This vector is the sole
piece of primal information exchanged in the forward direction: it
enters window $i{+}1$ as the parameter
\begin{equation}
\bm{\rho}_{i+1} \;\triangleq\; \big(\widehat{B}^{\,\text{prev}}_{i+1,j}\big)_{j\in\mathcal{B}},
\qquad
\bm{\rho}_{i+1} \,\leftarrow\, \bm{x}_i.
\end{equation}
The reverse-direction exchange carries dual information through Benders
cuts (Section~\ref{ssec:subproblem}).

% ---------------------------------------------------------------
\subsection{Window Sub-Problem}
\label{ssec:subproblem}

Each window $i$ solves a local OPF over $\mathcal{W}_i$, augmented by a
scalar epigraph variable $\theta_i\!\ge\!0$ representing the optimal
value of the \emph{downstream} problem (windows $i{+}1,\dots,P$) as a
function of the boundary SOC $\bm{x}_i$. Denoting by $\bm{z}_i$ the
local decision vector (line flows, nodal voltages, generator outputs,
battery powers and SOCs) and by $\mathcal{F}_i(\bm{\rho}_i)$ the
feasible set of window $i$ parameterized by the incoming SOC
$\bm{\rho}_i$, the sub-problem is
\begin{subequations}
\label{eq:subproblem}
\begin{align}
\big[\mathrm{SP}_i^{(k)}\big]\quad
V_i^{(k)}(\bm{\rho}_i) \;=\;
\min_{\bm{z}_i,\,\theta_i}\;
& c_i(\bm{z}_i)\;+\;\theta_i
\label{eq:sub_obj}\\[2pt]
\text{s.t.}\quad
& \bm{z}_i \,\in\, \mathcal{F}_i(\bm{\rho}_i),
\label{eq:sub_feas}\\[2pt]
& B_{w^{s}_i,j} \;=\; \rho_{i,j} \;+\; \eta_c P^{c}_{w^{s}_i,j}
   \;-\; \eta_d^{-1} P^{d}_{w^{s}_i,j},
   \;\;\forall j\!\in\!\mathcal{B},
\label{eq:sub_link}\\[2pt]
& \theta_i \;\ge\; \alpha_\ell \;+\; \bm{\beta}_\ell^{\!\top}\, \bm{x}_i,
   \quad \ell = 1,\dots,L_i^{(k)},
\label{eq:sub_cuts}\\[2pt]
& \theta_i \;\ge\; 0.
\label{eq:sub_thpos}
\end{align}
\end{subequations}
Here $c_i(\bm{z}_i)$ is the local stage cost (generation, voltage
penalty, storage degradation, etc.); \eqref{eq:sub_link} is the only
constraint where the incoming boundary $\bm{\rho}_i$ enters the
model; and \eqref{eq:sub_cuts} accumulates the $L_i^{(k)}$ Benders
optimality cuts received from window $i{+}1$ up to outer iteration $k$.
For the rightmost window we set $\theta_P\!\equiv\!0$, since no
downstream problem exists.

\paragraph*{What each window needs and provides}
Window $i$ requires two inputs at iteration $k$:
\textit{(i)} the latest forward state $\bm{\rho}_i^{(k)}$ produced by
window $i{-}1$, and
\textit{(ii)} the latest cut $(\alpha^{(k)},\bm{\beta}^{(k)})$
generated by window $i{+}1$. In return it provides three outputs:
\textit{(i)} its terminal SOC $\bm{x}_i^{(k)}$ for the forward pass,
\textit{(ii)} its objective value $V_i^{(k)}$, and
\textit{(iii)} the dual multiplier $\bm{\lambda}_i^{(k)}$ associated
with the boundary linking constraint~\eqref{eq:sub_link}. Crucially,
$\bm{\rho}_i$ enters the model \emph{only} through~\eqref{eq:sub_link},
so by sensitivity analysis
\begin{equation}
\label{eq:dual_sens}
\bm{\lambda}_i^{(k)} \;=\;
\frac{\partial V_i^{(k)}}{\partial \bm{\rho}_i}\bigg|_{\bm{\rho}_i^{(k)}},
\end{equation}
which is exactly the gradient information needed by Benders.

% ---------------------------------------------------------------
\subsection{Cut Generation: A Walk-Through of Window~$i$}
\label{ssec:cut_example}

To make the information flow concrete, consider an inner window
$i\in\{2,\dots,P{-}1\}$ at outer iteration $k$. The sequence of events
that involve window $i$ is the following.

\begin{enumerate}\itemsep2pt
\item[\textbf{S1.}]
\emph{Receive forward state.} Window $i$ reads
$\bm{\rho}_i^{(k)} = \bm{x}_{i-1}^{(k-1)}$ from the parent, i.e.\ the
terminal SOC delivered by window $i{-}1$ at the previous iteration.

\item[\textbf{S2.}]
\emph{Receive backward cut.} Window $i$ appends to its cut-pool the
optimality cut generated by window $i{+}1$ at iteration $k{-}1$:
\(
\theta_i \ge
\alpha_{i\leftarrow i{+}1}^{(k-1)} +
\bm{\beta}_{i\leftarrow i{+}1}^{(k-1)\,\top}\bm{x}_i.
\)

\item[\textbf{S3.}]
\emph{Solve.} It solves
$[\mathrm{SP}_i^{(k)}]$ to obtain $V_i^{(k)}$, the local primal
$\bm{z}_i^{(k)}$, the boundary dual
$\bm{\lambda}_i^{(k)}$ from~\eqref{eq:dual_sens}, and the new
terminal SOC $\bm{x}_i^{(k)} = (B_{w^{e}_i,j}^{(k)})_j$.

\item[\textbf{S4.}]
\emph{Emit forward state.} It ships $\bm{x}_i^{(k)}$ to window $i{+}1$,
which will be its $\bm{\rho}_{i+1}^{(k+1)}$ in S1 of the next iteration.

\item[\textbf{S5.}]
\emph{Emit backward cut.} Using the exact-supporting-hyperplane
construction, window $i$ builds a cut for window $i{-}1$'s epigraph
variable $\theta_{i-1}$:
\begin{equation}
\label{eq:cut_build}
\boxed{\;
\theta_{i-1} \;\ge\;
\underbrace{V_i^{(k)} \;-\; \bm{\lambda}_i^{(k)\,\top}\,\bm{\rho}_i^{(k)}}_{\displaystyle \alpha_{i-1\leftarrow i}^{(k)}}
\;+\;
\underbrace{\bm{\lambda}_i^{(k)}{}^{\top}}_{\displaystyle \bm{\beta}_{i-1\leftarrow i}^{(k)\,\top}}\,
\bm{x}_{i-1}\;.}
\end{equation}
Equation~\eqref{eq:cut_build} is the first-order Taylor expansion of the
downstream value function $V_i$ at the linearization point
$\bm{\rho}_i^{(k)}$, evaluated at the decision variable
$\bm{x}_{i-1}$ of the parent window. Because each $V_i$ is convex in
its boundary argument (the sub-problem is a convex QCP under the
SOCP relaxation of the branch-flow model), \eqref{eq:cut_build} is a
\emph{valid} under-estimator, and accumulating it yields an
increasingly tight piecewise-linear surrogate for $\theta_{i-1}$.
\end{enumerate}

Two boundary cases complete the picture. Window $1$ never receives a
forward state ($\bm{\rho}_1$ is fixed to the global initial SOC
$B_0$) and never emits a backward cut, but it consumes cuts from
window $2$. Window $P$ has $\theta_P\!\equiv\!0$, never receives a
backward cut, and emits both a forward state (used by no one) and a
backward cut for window $P{-}1$.

% ---------------------------------------------------------------
\subsection{Global Bounds and Convergence Criterion}
\label{ssec:bounds}

Because the windows are disjoint and exhaustive, at every outer
iteration $k$ the algorithm exposes two readily computable global
quantities:
\begin{align}
\underline{Z}^{(k)} \;&\triangleq\; V_1^{(k)}
\;=\; c_1(\bm{z}_1^{(k)}) + \theta_1^{(k)},
\label{eq:LB}\\[2pt]
\overline{Z}^{(k)} \;&\triangleq\; \sum_{i=1}^{P} c_i(\bm{z}_i^{(k)}).
\label{eq:UB}
\end{align}
The quantity $\underline{Z}^{(k)}$ is a \emph{global lower bound} on
the centralized optimum: $\theta_1$ is an under-estimator of the
true value-to-go of windows $2,\dots,P$, hence
$V_1^{(k)}\!\le\!Z^\star$ for every $k$ and is non-decreasing in $k$
under standard Benders monotonicity arguments \cite{XX}. Symmetrically,
$\overline{Z}^{(k)}$ is the \emph{exact} primal cost of the
window-stitched trajectory: with zero overlap each timestep is
optimized by exactly one window, so the stage costs sum to the true
primal objective with no double-counting and no need for an
overlap-stitching rule.

We declare convergence when either of the following holds:
\begin{align}
\textbf{(C1)}\;\;\;
\frac{\overline{Z}^{(k)} - \underline{Z}^{(k)}}
     {\max(1,|\overline{Z}^{(k)}|)} &\;<\; \varepsilon_{\text{gap}},
\label{eq:cv_gap}\\[2pt]
\textbf{(C2)}\;\;\;
\max\!\Big\{\!
\tfrac{|\underline{Z}^{(k)}{-}\underline{Z}^{(k-1)}|}{\max(1,|\underline{Z}^{(k)}|)},\;
\tfrac{|\underline{Z}^{(k-1)}{-}\underline{Z}^{(k-2)}|}{\max(1,|\underline{Z}^{(k)}|)}\!
\Big\} &\;<\; \varepsilon_{\text{lb}}.
\label{eq:cv_lb}
\end{align}
Criterion (C1) is the textbook Benders gap; (C2) is a two-step
LB-stability test that guards against the residual primal degeneracy
observed at window boundaries, which causes $\overline{Z}^{(k)}$ to
oscillate by a few cuts even after the lower bound has plateaued.
Empirically, (C2) triggers first on networks with many co-located
batteries.

% ---------------------------------------------------------------
\subsection{The Parallel DDDP--OTD Algorithm}
\label{ssec:algorithm}

The full procedure is summarized in Algorithm~\ref{alg:dddp_otd}. The
key structural property is that, within an outer iteration, all $P$
sub-problems $[\mathrm{SP}_i^{(k)}]$ are \emph{simultaneously} solvable:
each one consumes only forward states and backward cuts from
\emph{previous} iterations (lines~\ref{algln:ship}--\ref{algln:recv}),
so no in-iteration synchronization between windows is required. The
parent process therefore dispatches all $P$ solves in parallel, blocks
on the slowest one, and only then performs the (cheap) cut-assembly
and state-pass bookkeeping. Communication is reduced to two small
dictionaries per window per iteration: a forward
$\bm{x}_i\!\in\!\mathbb{R}^{|\mathcal{B}|}$ and a backward
$(\alpha,\bm{\beta})\!\in\!\mathbb{R}^{1+|\mathcal{B}|}$.

\begin{algorithm}[t]
\caption{Parallel multi-cut DDDP--OTD}
\label{alg:dddp_otd}
\begin{algorithmic}[1]
\Require horizon $T$, partitions $P$, initial SOC $\bm{B}_0$,
         tolerances $\varepsilon_{\text{gap}},\varepsilon_{\text{lb}}$,
         max iterations $K$.
\State Build disjoint windows $\{\mathcal{W}_i\}_{i=1}^{P}$ with
       $w^{s}_1{=}1$, $w^{e}_P{=}T$, $w^{s}_{i+1}{=}w^{e}_i{+}1$.
\State Spawn $P$ persistent worker processes; each builds its local
       model $[\mathrm{SP}_i^{(0)}]$ \emph{once} with empty cut-pool.
\State Initialize $\bm{\rho}_i^{(1)} \leftarrow \bm{B}_0$ and
       pending cut-pool $\mathcal{C}_i \leftarrow \emptyset$
       for all $i\!\in\!\{1,\dots,P\}$.
\For{$k = 1, 2, \dots, K$}
   \State \textbf{(Parallel solve)}\;
          \textbf{parfor} $i = 1, \dots, P$ \textbf{do}
                          \label{algln:par_begin}
   \State \quad append cuts $\mathcal{C}_i$ to $[\mathrm{SP}_i^{(k)}]$
          and clear $\mathcal{C}_i$
                          \label{algln:recv}
   \State \quad solve $[\mathrm{SP}_i^{(k)}]$ at $\bm{\rho}_i^{(k)}$,
          obtaining $V_i^{(k)},\bm{x}_i^{(k)},\bm{\lambda}_i^{(k)}$
   \State \textbf{end parfor}
                          \label{algln:par_end}
   \State \textbf{(Cut assembly)} \For{$i = 2, \dots, P$}
       \State $\alpha \leftarrow V_i^{(k)} -
              \bm{\lambda}_i^{(k)\top}\bm{\rho}_i^{(k)}$;\;\;
              $\bm{\beta} \leftarrow \bm{\lambda}_i^{(k)}$
       \State $\mathcal{C}_{i-1} \leftarrow \mathcal{C}_{i-1}\cup\{(\alpha,\bm{\beta})\}$
                          \label{algln:cut_emit}
   \EndFor
   \State \textbf{(State pass)} for $i=1,\dots,P{-}1$:
          $\bm{\rho}_{i+1}^{(k+1)} \leftarrow \bm{x}_i^{(k)}$
                          \label{algln:ship}
   \State $\underline{Z}^{(k)} \leftarrow V_1^{(k)}$;\;\;
          $\overline{Z}^{(k)} \leftarrow \sum_{i=1}^{P} c_i(\bm{z}_i^{(k)})$
   \If{(C1) in \eqref{eq:cv_gap} or (C2) in \eqref{eq:cv_lb} holds}
       \State \textbf{break}
   \EndIf
\EndFor
\State Shut down workers; merge per-window primals into the global trajectory.
\State \Return $\{\bm{z}_i^{(k)}\}_{i=1}^P,\;
                \underline{Z}^{(k)},\;\overline{Z}^{(k)}$.
\end{algorithmic}
\end{algorithm}

% ---------------------------------------------------------------
\subsection{Discussion}
\label{ssec:discussion}

Three properties distinguish DDDP--OTD from its overlapping-Schwarz
counterpart and from textbook SDDP. First, the windows are
\emph{strictly disjoint}, which makes the upper bound $\overline{Z}$
exact rather than an inner approximation; no overlap weight has to be
tuned. Second, the multi-cut multi-window dual-passing structure
generalizes the two-stage Benders idea to a chain of $P$ stages while
remaining \emph{fully parallel} per iteration --- unlike SDDP, which
sweeps forward then backward sequentially. Third, the entire boundary
coupling is captured by the dual of a single equality constraint
(\eqref{eq:sub_link}), which makes the implementation agnostic to the
particular OPF formulation: any convex relaxation of the branch-flow
or bus-injection model can be dropped into $\mathcal{F}_i$ without
altering the cut machinery.

% =====================================================================
%  End of section
% =====================================================================
